% Publication figure captions; all 1,500 trials passed the final audit.

% simulation-design.pdf
\caption{\textbf{Simulated clade origins and their survival under study sampling.}
Examples use the first predetermined trial at the 50\% node-selection level
in the \textbf{a,} 0--80, \textbf{b,} 80--160, and \textbf{c,} 160--280~Ma epochs,
with 270, 89, and 9 binary characters, respectively. Each marked internal node
defines one character whose descendants have state~1. Orange marks characters
remaining variable among the 77 study taxa; dark grey marks characters whose
state-1 descendants are absent. All trees show the 765-tip phylogeny, but
origins are restricted to branches within Apocrita. Preservation of a variable
character does not imply preservation of its original branching node.}

% recovery-summary.pdf
\caption{\textbf{Recovery of simulated phenomic change under the study's taxon sampling.}
\textbf{a,} The fraction of characters remaining variable quantifies direct
loss of the imposed signal when reducing the 765-tip tree to the 77 study taxa.
\textbf{b,} The study/full ratio of combined rates in the target epoch captures
the joint effects of character loss, temporal redistribution, and the different
lineage-time denominators of the two trees; it is not a direct estimate of
bias in the fitted transition matrix.
\textbf{c,} The fraction of branch-averaged reconstructed study-tree change
assigned to the selected node-age epoch measures temporal localization.
\textbf{d,} The fraction assigned to each character's own origin branch
distinguishes branch assignment from temporal localization on long branches.
Points and lines show means over 100 trials per condition; bars show empirical
central 95\% ranges across trials. Colours distinguish the three epochs.
The last two panels use trials with surviving variable characters.
Dotted lines mark complete retention or localization, or an epoch-rate ratio
of one. Requested percentages are rounded to integer node counts as described
in the Methods.}

% marginal-rates-all.pdf, three pages
\caption{\textbf{Recovery of temporal increases in phenomic change after taxon subsampling.}
The three pages show origins at nodes aged 0--80, 80--160 and 160--280~Ma.
Panels show 5\%, 10\%, 30\%, 50\% and 80\% node selection from top to bottom;
character counts refer to the initial simulated datasets.
Black solid lines show the 765-taxon tree and orange dashed lines the
77-taxon study tree. Rates are summed across characters and averaged over
lineages present at each age. Lines show the average of 100 trial means; coloured bands enclose
their central 95\% range.
Grey backgrounds mark the simulated node-age interval. Characters made
constant after pruning are omitted from study-tree reconstruction.
Branch averaging can spread reconstructed change beyond the selected interval.}

% marginal-rates-matched.pdf, three pages
\caption{\textbf{Temporal rate reconstruction using the same surviving characters.}
Both trees use exactly the characters remaining variable after pruning to
the 77 study taxa.
The three pages show origins at nodes aged 0--80, 80--160 and 160--280~Ma;
panels show 5\%, 10\%, 30\%, 50\% and 80\% node selection from top to bottom. Panel counts refer to initial
characters, whereas the number retained varies among trials.
Black solid lines show the 765-taxon tree and orange dashed lines the
77-taxon tree. Rates are summed across characters and averaged over
lineages present at each age. Lines show the average of 100 trial means; coloured bands enclose
their central 95\% range. Grey backgrounds mark the simulated node-age interval.}
